\section{Database Design}

\subsection{Storage Strategy}

Persistence is divided across four stores rather than one. The division is deliberate:
the three kinds of state the system holds have incompatible access patterns, and forcing
them into a single store would compromise all three.

\begin{table}[htbp]
\centering
\small
\begin{tabularx}{\textwidth}{@{}L{3.1cm}L{3.2cm}XL{1.9cm}@{}}
\toprule
\textbf{Store} & \textbf{Technology} & \textbf{Contents} & \textbf{Unit} \\
\midrule
Application database & Supabase Postgres &
Accounts, projects, the recording registry, conversations and their messages. Normalised,
transactional, and subject to row-level security. & Row \\
Object storage & Supabase Storage, two buckets &
Video bytes as uploaded, the canonical copy the pipeline reads, and one poster frame per
recording. & Object \\
Analysis record store & JSON documents on local disk &
Everything the analyzers and aggregators produced for one recording, nested under its
chunks. & Document \\
Vector store & Qdrant, embedded &
One point per chunk per analyzer, carrying five embeddings and a filterable payload. &
Point \\
\bottomrule
\end{tabularx}
\end{table}

The analysis output is held as documents rather than as relations because its schema is
defined by the analyzer registry, not by the database: each analyzer declares its own
output shape, and a chunk carries one key per analyzer that ran on it. Modelling that
relationally would make adding an analyzer a schema migration, which directly contradicts
NFR-8. Conversely, application state is relational because it requires foreign keys,
cascading deletes and per-row access control, none of which a document store provides as
cheaply. Retrieval is separated again because approximate nearest-neighbour search over
several named vector spaces is not something either of the other two can serve.

\subsection{Application Database}

Figure~\ref{fig:erd} gives the logical structure --- which relations exist and how they
are related --- and Figure~\ref{fig:schema} the physical schema as it stands in Postgres,
with every column, its type and whether it is nullable.

\subsubsection*{Entity relationships}

\begin{figure}[htbp]
\centering
\begin{tikzpicture}

  \node[extbox] (users) {\textbf{auth.users}\\[2pt]\scriptsize Supabase Auth};
  \node[entity, below=11mm of users] (projects)
    {\textbf{projects}\\[2pt]\scriptsize id, user\_id, name};

  \node[entity, below left=12mm and 5mm of projects] (video)
    {\textbf{video\_core}\\[2pt]\scriptsize id, project\_id, user\_id,\\
     \scriptsize core\_video\_id, status};
  \node[entity, below right=12mm and 5mm of projects] (conv)
    {\textbf{conversations}\\[2pt]\scriptsize id, project\_id, user\_id};
  \node[entity, below=11mm of conv] (msg)
    {\textbf{messages}\\[2pt]\scriptsize id, conversation\_id, role, parts};

  \draw[rel] (users) -- node[card, right=1mm] {$1 : n$} (projects);
  \draw[rel] (projects.south) -- ++(0,-4mm) -| node[card, pos=0.3, above] {$1 : n$} (video.north);
  \draw[rel] (projects.south) -- ++(0,-4mm) -| node[card, pos=0.3, above] {$1 : n$} (conv.north);
  \draw[rel] (conv) -- node[card, right=1mm] {$1 : n$} (msg);

\end{tikzpicture}
\caption{Application database. Every relationship cascades on delete, so removing a
project removes its recordings, conversations and messages with it.}
\label{fig:erd}
\end{figure}

\subsubsection*{Physical schema}

\optfigure{falcon-vqa-schema.png}{Application database as deployed: the four live
relations with their columns, types and nullability, and their foreign keys into
\texttt{auth.users}. Filled markers denote non-nullable columns.}{0.92}{fig:schema}

\subsubsection*{\texttt{video\_core}}

This relation is the application's view of a recording analysed by the core service. It
is the only place where the two halves of the system meet, and three of its design
decisions are load-bearing.

\begin{table}[htbp]
\centering
\small
\begin{tabularx}{\textwidth}{@{}L{3.2cm}L{2.5cm}X@{}}
\toprule
\textbf{Column} & \textbf{Type} & \textbf{Purpose} \\
\midrule
\texttt{id} & \texttt{uuid} PK & Application identity, independent of the analysis service. \\
\texttt{project\_id} & \texttt{uuid} FK & Owning project; cascades on delete. \\
\texttt{user\_id} & \texttt{uuid} FK & Owning account; the column row-level security tests. \\
\texttt{title} & \texttt{text} & Display name, editable without affecting analysis. \\
\texttt{source\_type} & \texttt{text} & \texttt{upload} or \texttt{url}. \\
\texttt{storage\_path} & \texttt{text} & Object in the application project's bucket. \\
\texttt{source\_url} & \texttt{text} & The URL the analysis service was handed. \\
\texttt{core\_video\_id} & \texttt{text} & The content hash the analysis service identifies this recording by. Null until the bytes have been fetched. \\
\texttt{playback\_url} & \texttt{text} & The public MP4 in the analysis project's bucket; this is what actually plays. \\
\texttt{poster\_url} & \texttt{text} & Single frame written at ingest, for the workspace grid. \\
\texttt{duration}, \texttt{size\_bytes} & \texttt{numeric}, \texttt{bigint} & Reported back by the analysis service. \\
\texttt{status} & \texttt{text} & \texttt{pending}, \texttt{uploading}, \texttt{queued}, \texttt{analyzing}, \texttt{ready} or \texttt{failed}. \\
\texttt{job\_id} & \texttt{text} & The background job being polled. \\
\texttt{stage} & \texttt{text} & \texttt{fetching}, \texttt{chunking}, \texttt{analyzing}, \texttt{indexing}, \texttt{aggregating} or \texttt{complete}. \\
\texttt{progress} & \texttt{jsonb} & The job's last progress detail, rendered directly by the interface. \\
\texttt{error} & \texttt{text} & Failure reason, retained so the row explains itself. \\
\texttt{ingest\_config} & \texttt{jsonb} & What the pipeline was asked for: analyzers, mode, preset, weights, interval and duration bounds. \\
\texttt{analyzers} & \texttt{text[]} & What it actually produced. \\
\texttt{aggregates} & \texttt{text[]} & Which video-level results exist. \\
\texttt{chunk\_config} & \texttt{text} & The chunking scheme, e.g. \texttt{audio\_video:5-20}. \\
\texttt{chunk\_count} & \texttt{int} & Number of chunks produced. \\
\texttt{created\_at}, \texttt{updated\_at} & \texttt{timestamptz} & \texttt{updated\_at} is maintained by trigger. \\
\bottomrule
\end{tabularx}
\end{table}

\begin{enumerate}[label=\textbf{D-\arabic*.},leftmargin=3.0em]
  \item \textbf{Two URLs are kept, not one.} The system spans two Supabase projects: this
        one holds the row and the object the browser uploaded, while the analysis
        service's own project holds the canonical copy that plays.
        \texttt{source\_url} and \texttt{playback\_url} therefore point at different
        objects and both are retained.
  \item \textbf{Uniqueness is per project, never global.} Because
        \texttt{core\_video\_id} is the content hash, the same file added to two
        projects yields the same value. The unique index is declared over
        \texttt{(project\_id, core\_video\_id)} and is partial --- restricted to rows
        where the column is non-null, since the identifier does not exist until the
        bytes have been downloaded. Deleting one project's row must therefore not
        delete the analysis service's copy while another row still references it.
  \item \textbf{Produced analyzers are an array, not a document.} Both the interface and
        the agent evaluate \emph{``does this recording have people analysis?''} as a
        predicate, which an array answers directly and an opaque JSON document does not.
        The requested configuration, by contrast, is \texttt{jsonb}, because it is
        replayed verbatim on re-index rather than queried.
\end{enumerate}

\subsubsection*{Remaining relations}

\begin{table}[htbp]
\centering
\small
\begin{tabularx}{\textwidth}{@{}L{2.6cm}X@{}}
\toprule
\textbf{Relation} & \textbf{Columns and notes} \\
\midrule
\texttt{projects} & \texttt{id}, \texttt{user\_id}, \texttt{name}, \texttt{description}, timestamps. The grouping every other relation hangs from. \\
\texttt{conversations} & \texttt{id}, \texttt{user\_id}, \texttt{project\_id}, \texttt{title}, \texttt{lastContext}, timestamps. A conversation is optionally attached to a project; the attachment is what scopes the agent's tools. \\
\texttt{messages} & \texttt{id}, \texttt{conversation\_id}, \texttt{role}, \texttt{parts}, \texttt{metadata}, \texttt{created\_at}. \texttt{parts} is \texttt{jsonb} because a turn is a sequence of heterogeneous parts --- text, tool calls, tool results and artifacts --- whose shapes are set by the agent runtime. \\
\texttt{videos} & Superseded. The relation predates the current analysis service and is retained but no longer read by anything; \texttt{video\_core} replaced it additively so that no migration had to drop data. \\
\bottomrule
\end{tabularx}
\end{table}

\subsubsection*{Indexes, integrity and access control}

\begin{itemize}
  \item \textbf{Indexes.} \texttt{video\_core} is indexed on
        \texttt{(project\_id, created\_at DESC)} for the workspace listing, and partially
        on \texttt{job\_id} where it is non-null, which is the lookup that reconciliation
        performs. \texttt{conversations} is indexed on \texttt{project\_id} and
        \texttt{messages} on \texttt{(conversation\_id, created\_at)}, the order a
        conversation is replayed in.
  \item \textbf{Integrity.} Every foreign key cascades on delete, so a removed account or
        project leaves nothing orphaned. A shared trigger function maintains
        \texttt{updated\_at} on each relation that carries it.
  \item \textbf{Access control.} Row-level security is enabled on every relation.
        \texttt{projects}, \texttt{conversations} and \texttt{video\_core} test
        \texttt{auth.uid() = user\_id} directly. \texttt{messages} carries no owner
        column, so its policy tests ownership of the parent conversation through an
        \texttt{EXISTS} subquery --- the row is reachable only if the conversation it
        belongs to is. This is the database half of NFR-9; the other half is the
        project-scoping function through which the agent's tools resolve recording
        identifiers.
\end{itemize}

\subsection{Object Storage}

Two public buckets are used, one in each Supabase project.

\begin{table}[htbp]
\centering
\small
\begin{tabularx}{\textwidth}{@{}L{2.6cm}L{2.7cm}X@{}}
\toprule
\textbf{Bucket} & \textbf{Project} & \textbf{Contents and policy} \\
\midrule
\texttt{project-assets} & Application &
What the browser uploads, laid out per user and project. Public read; insert restricted
to authenticated accounts; delete restricted to the object's owner. \\
\texttt{videos} & Analysis service &
The canonical copy the pipeline reads, stored at
\texttt{\{video\_id\}/\{filename\}}, together with
\texttt{\{video\_id\}/poster.jpg}. Public read, so a playback URL is permanent and
unsigned. \\
\bottomrule
\end{tabularx}
\end{table}

Because the object path begins with the content hash, re-ingesting identical bytes
resolves to an object that already exists and the upload is skipped entirely. The
service checks for the object before writing rather than relying on overwrite semantics,
so a repeated ingest of a large recording costs nothing.

\subsection{Analysis Record Store}

One JSON document is written per recording, named
\texttt{\{sanitised-filename\}\_\_\{video\_id[:8]\}.json}. The filename carries a
human-readable stem so that a directory of records can be inspected directly, and the
truncated content hash so that two recordings sharing a name cannot collide. On
re-ingest, any existing document for the same identifier is removed before the new one
is written, so a stale record can never sit alongside a current one under a different
name.

\begin{table}[htbp]
\centering
\small
\begin{tabularx}{\textwidth}{@{}L{3.4cm}L{2.2cm}X@{}}
\toprule
\textbf{Field} & \textbf{Type} & \textbf{Contents} \\
\midrule
\texttt{video\_id} & string & Content hash; the identity every other store joins on. \\
\texttt{video\_url} & string & Public URL of the canonical copy. \\
\texttt{storage\_path} & string & Object path within the analysis bucket. \\
\texttt{source\_url} & string & Where the recording was fetched from. \\
\texttt{filename} & string & Original name, preserved for display. \\
\texttt{preset}, \texttt{params} & string, object & The chunking scheme and its duration bounds or weights. \\
\texttt{chunk\_config} & string & The composite key identifying that scheme. \\
\texttt{duration}, \texttt{size\_bytes} & number & Measured at ingest. \\
\texttt{poster\_url} & string & The extracted poster frame. \\
\texttt{analyzers} & array & Which analyzers ran. \\
\texttt{chunks} & array & One entry per chunk; see below. \\
\texttt{aggregates} & object & One key per aggregator, holding its full result. \\
\texttt{aggregates\_analyzers} & array & The analyzer set the aggregates were computed from, so a stale cache is detectable. \\
\bottomrule
\end{tabularx}
\end{table}

Each entry of \texttt{chunks} carries \texttt{id}, \texttt{start} and \texttt{end},
followed by \emph{one key per analyzer that produced output for that chunk}, whose value
is that analyzer's structured result --- \texttt{default\_video} holding the scene
schema, \texttt{people} holding per-person records with their box locations,
\texttt{ocr} holding the read texts, and so on. An analyzer that was gated out for a
chunk simply contributes no key, so absence is represented naturally rather than as a
null. This is precisely the shape a relational schema cannot accommodate without a
migration per analyzer.

Two conventions apply throughout. Keys prefixed with an underscore ---
\texttt{\_frames}, \texttt{\_detector\_labels} --- are provenance retained for debugging:
they record which frames a description was derived from and what the cheap detector
believed, and they are never embedded, never copied into a vector payload and never
returned by the API. And no local filesystem path is ever persisted; the cached file is
re-derived from \texttt{video\_id} on demand and re-downloaded if the cache is cold, so
a record remains correct after the data directory moves or the cache is wiped.

\subsection{Vector Store}

\subsubsection*{Collection and point structure}

A single Qdrant collection, \texttt{chunks}, holds one point per
\emph{(recording, analyzer, chunk span, chunk configuration)}. Each point carries all
five named vectors, each 384-dimensional and compared by cosine distance. Making the
point --- rather than the vector --- the unit of storage means a limit of ten returns ten
chunks, the payload is stored once, and a filter is evaluated once.

\begin{table}[htbp]
\centering
\small
\begin{tabularx}{\textwidth}{@{}L{2.5cm}X@{}}
\toprule
\textbf{Named vector} & \textbf{Embedded text} \\
\midrule
\texttt{combined}    & The whole flattened record. The default search space. \\
\texttt{description} & The prose description alone. \\
\texttt{people}      & Role, clothing and appearance, phrased as clauses. \\
\texttt{actions}     & What each subject did, or what each object was used for. \\
\texttt{objects}     & Object names, or the literal strings read from screen. \\
\bottomrule
\end{tabularx}
\end{table}

A vector space is omitted rather than written empty, so a chunk in which no person was
seen is not a candidate in a people-scoped search at all, instead of matching against
empty text.

\subsubsection*{Point identity}

The point identifier is a UUID derived deterministically from

\begin{center}
\texttt{sha1(video\_id | analyzer\_id | start | end | chunk\_config)}
\end{center}

so that re-ingesting the same chunk upserts in place rather than duplicating. The key is
built from the \emph{time span} rather than a positional index, because chunk 12 of one
chunking run is a different moment from chunk 12 of another and a positional identifier
would silently rebind a vector to the wrong timeframe. Custom chunking weights are
hashed into \texttt{chunk\_config} for the same reason: without it, two different
weightings sharing the same duration bounds would collide and the second ingest would
overwrite the first.

\subsubsection*{Payload and filters}

The payload carries everything required to rank, filter and cite a moment, so retrieval
never joins back to the record store: the recording identifier and playback URL, the
analyzer, chunk configuration and chunk identifier, the time span, the combined text and
description, the setting, the people, objects, actions and tags, the speakers and their
turns, the texts read from screen, the object detections, the full per-person records,
and the person count.

\begin{table}[htbp]
\centering
\small
\begin{tabularx}{\textwidth}{@{}L{2.5cm}L{2.4cm}L{1.9cm}X@{}}
\toprule
\textbf{Filter} & \textbf{Payload field} & \textbf{Match} & \textbf{Type} \\
\midrule
\texttt{video\_ids}    & \texttt{video\_id}     & any   & \texttt{list[str]} \\
\texttt{chunk\_ids}    & \texttt{chunk\_id}     & any   & \texttt{list[int]} \\
\texttt{analyzer\_ids} & \texttt{extractor\_id} & any   & \texttt{list[str]} \\
\texttt{chunk\_config} & \texttt{chunk\_config} & exact & \texttt{str} \\
\texttt{objects}       & \texttt{objects}       & any   & \texttt{list[str]} \\
\texttt{tags}          & \texttt{tags}          & any   & \texttt{list[str]} \\
\texttt{speakers}      & \texttt{speakers}      & any   & \texttt{list[str]} \\
\texttt{people}        & \texttt{people}        & any   & \texttt{list[str]} \\
\texttt{min\_people}   & \texttt{people\_count} & $\geq$ & \texttt{int} \\
\texttt{max\_people}   & \texttt{people\_count} & $\leq$ & \texttt{int} \\
\texttt{after}         & \texttt{start}         & $\geq$ & \texttt{float} \\
\texttt{before}        & \texttt{end}           & $\leq$ & \texttt{float} \\
\bottomrule
\end{tabularx}
\end{table}

The filter set is declared in one place and the API passes a flat dictionary through to
it, so a filter added to the declaration is reachable from the API immediately; the
previous hand-plumbed parameter list left six store filters silently unreachable. An
unrecognised filter name raises an error carrying the nearest valid name, because a
filter that is silently ignored returns plausible but wrong results, which is worse than
a failure.

The same twelve fields are declared as payload indexes. The person count is indexed
specifically because counting is what embeddings cannot do: a query for \emph{``mostly
empty store''} retrieved the busiest chunk in the recording, since a description listing
five people reads much the same to a vector as one listing two. In embedded mode Qdrant
ignores payload indexes and filtering scans instead --- correct, and acceptable at this
scale. The index declarations are retained so that pointing the same code at a Qdrant
server gains real indexes with no change.

Deletion is scoped rather than wholesale. Removing vectors takes an optional chunk
configuration and analyzer identifier, because re-running one analyzer must delete only
that analyzer's vectors; deleting by recording alone would destroy every other
analyzer's work, so adding one analyzer to an existing recording would silently discard
the rest.

\subsection{Cross-Store Identity}

The content hash is the single join key across all four stores. It is computed while the
bytes are being read, which is why the download must complete before the object's own
name exists.

\begin{figure}[htbp]
\centering
\begin{tikzpicture}

  \node[storebox] (pg)
    {\textbf{Postgres}\\[2pt]\scriptsize \texttt{video\_core}\\
     \scriptsize \texttt{core\_video\_id}};
  \node[storebox, right=8mm of pg] (obj)
    {\textbf{Object storage}\\[2pt]\scriptsize \texttt{videos} bucket\\
     \scriptsize \texttt{\{hash\}/file.mp4}};
  \node[storebox, right=8mm of obj] (rec)
    {\textbf{Record store}\\[2pt]\scriptsize \texttt{\{hash\}.json}\\
     \scriptsize \texttt{video\_id}};
  \node[storebox, right=8mm of rec] (qd)
    {\textbf{Qdrant}\\[2pt]\scriptsize \texttt{chunks} payload\\
     \scriptsize \texttt{video\_id}};

  \draw[keyflow] (pg) -- (obj);
  \draw[keyflow] (obj) -- (rec);
  \draw[keyflow] (rec) -- (qd);

  \node[card, below=4mm of obj.south] (note)
    {joined by one content hash, \texttt{sha1(bytes)[:16]}};

\end{tikzpicture}
\caption{The content hash of the video bytes is the identity shared by every store, which
is what makes ingestion idempotent (NFR-10).}
\label{fig:keys}
\end{figure}

Durability differs by store, and the difference is intentional:

\begin{itemize}
  \item The \textbf{record store} is authoritative. It holds the analyzer and aggregator
        output, which is the expensive artifact, and everything else can be rebuilt from
        it.
  \item The \textbf{vector store} is derived. Re-indexing rebuilds it from the records
        at the cost of local embedding alone, with no API spend and no re-analysis.
  \item The \textbf{media cache} is regenerable. Deleting it costs a re-download and
        never data, which is why no local path is persisted anywhere.
  \item \textbf{Job state} is deliberately in memory and does not survive a restart. The
        application detects the loss when a poll returns nothing, and either recovers
        the recording from its persisted record or fails the row explicitly so that
        re-indexing becomes the route forward.
\end{itemize}
